% 【本文件写什么】impl 实现层：dummy
% - 定位：占位/链接桩实现，保证默认 feature 可编译。
% - crate 路径：os/components/wateros-fs/fs-devfs/devfs-impl/impl-dummy/src/
% - 如何实现对应 api-v0 trait：关键类型、状态机、数据结构
% - 算法或硬件交互流程，可用伪代码/流程图
% - 本 impl 启用的 Cargo [features] 与 arch feature，impl-riscv64 等
% - 与其它 impl 变体的差异、切换 feature 时的影响
% - 性能/内存假设与已知限制
% 【事实来源】
% - 上述 crate 源码与 Cargo.toml
% - os/feature-tree.txt 中指向本 impl 的 feature 链

\subsubsection{impl-dummy}

\code{DummyDevFsManager}实现\code{DevFsManager}全部方法但无真实设备：\code{refresh}为空操作，\code{list\_nodes}返回空表，\code{lookup\_*}与\code{register\_*}返回\code{NotFound}/\code{Unsupported}，\code{default\_root\_block\_path}恒为\code{None}。

用于无驱动或最小链接配置；\code{fs-devfs}在同时启用\code{impl-kernel}与\code{impl-dummy}时优先选择\code{impl-kernel}。启用dummy后\code{fs::init()}无法获得根块设备，probe链在devfs阶段即终止。
